% id: 132-320
% book: 개념원리 대수 · 14 일반각
% section: 개념원리 익히기
% figure: none
% answer: 풀이 참조 (⑴ 시초선~$\pt{OX}$에서 양의 방향으로 $45^\circ$ 회전한 동경~$\pt{OP}$ ⑵ 음의 방향으로 $250^\circ$ 회전한 동경~$\pt{OP}$ ⑶ 양의 방향으로 $630^\circ$(한 바퀴와 $270^\circ$) 회전한 동경~$\pt{OP}$ ⑷ 음의 방향으로 $750^\circ$(두 바퀴와 $30^\circ$) 회전한 동경~$\pt{OP}$의 그림)
% answer_source: 답지
% variant_level: 0 (원본 전사)
% 이 파일은 build.mjs 가 items.json 에서 만든다. 고칠 때는 items.json 을 고친다.
\smash{\raisebox{1.5mm}{\dmkichul{개념원리 대수 132쪽 320번}}}
다음 각을 나타내는 시초선~$\pt{OX}$와 동경~$\pt{OP}$의 위치를 그림으로 나타내시오.
\dmsubprob{1} $45^\circ$
\dmsubprob{2} $-250^\circ$
\dmsubprob{3} $630^\circ$
\dmsubprob{4} $-750^\circ$
